\documentclass[a4paper,11pt]{article}
\usepackage[utf8]{inputenc}
\usepackage[T1]{fontenc}
\usepackage[portuguese]{babel}
\usepackage{amsmath,amssymb}
\usepackage[margin=2.5cm]{geometry}

% Um passo da prova: \passo{profundidade}{(n)}{texto}{conjunto}
\newcommand{\passo}[4]{%
  \par{\leftskip=\dimexpr 2em*#1\relax
  \noindent\makebox[2.5em][l]{#2}#3\par
  \if\relax\detokenize{#4}\relax\else\noindent\hspace*{2.5em}#4\par\fi}%
  \smallskip}

% Pares resolvidos de uma eliminação: (cláusula com x) · (cláusula com ¬x) → resolvente
\newenvironment{pares}{\begin{itemize}\setlength{\itemsep}{0pt}\small}{\end{itemize}\smallskip}

\begin{document}

\section*{Algoritmo DP (Davis-Putnam, 1960)}

Conjunto de cláusulas: $S = \{\neg p \lor \neg q,\allowbreak\ \neg p \lor q,\allowbreak\ p \lor \neg q,\allowbreak\ p \lor q\}$

\subsection*{Sequência de conjuntos}

O DP não ramifica nem volta atrás: em cada passo troca $S$ por um conjunto que é satisfazível se e só se $S$ o for. Sem cláusulas unitárias, elimina a primeira variável $x$ por resolução: as cláusulas com $x$ e com $\neg x$ dão lugar aos resolventes de cada par.

\passo{0}{(1)}{Conjunto inicial $S$:}{$\{\neg p \lor \neg q,\allowbreak\ \neg p \lor q,\allowbreak\ p \lor \neg q,\allowbreak\ p \lor q\}$}
\passo{0}{(2)}{Eliminação de $p$ por resolução:}{$\{\neg q,\allowbreak\ q\}$}
\noindent\hspace*{2.5em}{\small Pares resolvidos (cláusula com $p$ $\cdot$ cláusula com $\neg p$ $\to$ resolvente):}
\begin{pares}
\item $p \lor \neg q \;\cdot\; \neg p \lor \neg q \;\to\; \neg q$
\item $p \lor \neg q \;\cdot\; \neg p \lor q \;\to\; \neg q \lor q$ (tautologia: descartada)
\item $p \lor q \;\cdot\; \neg p \lor \neg q \;\to\; q \lor \neg q$ (tautologia: descartada)
\item $p \lor q \;\cdot\; \neg p \lor q \;\to\; q$
\end{pares}
\passo{0}{(3)}{Propagação de $\neg q$:}{$\{F\}$}
\passo{0}{}{$S$ contém $F$ $\Rightarrow$ insatisfazível}{}
\passo{0}{}{\textbf{Resultado: $S$ é insatisfazível.}}{}

\vfill
{\footnotesize Gerado com Learn Davis-Putnam SAT, desenvolvido por Miguel Rocha (miguelrocha.dev).}

\end{document}
